\documentclass{article}
\usepackage[UTF8]{ctex}
\usepackage{amsmath,amssymb,amsfonts}
\usepackage{enumitem}
\usepackage{geometry}
\geometry{a4paper, margin=1in}
% === 错题本排版优化（不参与题目标识） ===
\setlength{\parindent}{0pt}
\setlength{\parskip}{0.35em plus 0.10em minus 0.08em}
\setlength{\emergencystretch}{3em}
\setlist{itemsep=0.25em, topsep=0.40em, parsep=0pt}
\allowdisplaybreaks[2]
% === 错题本排版优化结束 ===
\begin{document}
\section*{西北工业大学2026年数学分析真题}
\section*{183．西北工业大学2026年数学分析试卷}

\subsection{题目1}
一．（20 分）用极限的严格数学定义证明：

\begin{itemize}
  \item[（1）] \(\lim _{x \rightarrow+\infty} \frac{x^{2}+x-1}{x^{2}}=1\) ．
\end{itemize}

\subsection{题目1}
原题边界“（2）”：\(\lim _{n \rightarrow \infty} \sqrt[n]{a}=1\) ，其中 \(a>1\) ．

\subsection{题目2}
原题边界“二．”：（20 分）解答如下问题：

\begin{itemize}
\item[（1）] 求极限： \(\lim _{x \rightarrow+\infty}\left(\sqrt{x^{4}+1}-\left(x^{2}-1\right) e^{\frac{1}{x}}\right)\) ．
\end{itemize}

\subsection{题目2}
原题边界“（2）”：求幂级数 \(\sum_{n=1}^{\infty} \frac{x^{n}}{n(n+1) 2^{n}}\) 的收敛域与和函数．

\subsection{题目2}
三．（15 分）使用闭区间套定理证明聚点定理：实数轴上任一有界无限点集 \(S\) 必有一个聚点．

\subsection{题目3}
四．（15 分）证明：若 \(f(x)\) 在 \([a, b]\) 上黎曼可积，则 \(\lim _{p \rightarrow+\infty} \int_{a}^{b} f(x) \cos p x \mathrm{~d} x=0\) ．

\subsection{题目4}
五．（15 分）利用 \(\int_{0}^{+\infty} e^{-x^{2}} \mathrm{~d} x=\frac{\sqrt{\pi}}{2}\) ，计算 \(\int_{0}^{+\infty} \frac{e^{-a^{2} x^{2}}-e^{-b^{2} x^{2}}}{x^{2}} \mathrm{~d} x\) ，其中 \(b>a>0\) ．

\subsection{题目5}
六。（15 分）函数 \(z=z(x, y)\) 二阶偏导连续，试以 \(u, v\) 作为新的自变量变换方程

\[
x^{2} \frac{\partial^{2} z}{\partial x^{2}}-y^{2} \frac{\partial^{2} z}{\partial y^{2}}=0
\]

其中 \(u=x y, v=\frac{x}{y}\) ．

\subsection{题目6}
七．（15 分）证明：\(f(x)=\sqrt{x}\) 在 \([0,+\infty)\) 上一致连续．

\subsection{题目7}
八．（15 分）应用高斯公式计算三重积分

\[
\iiint_{V}(x y+y z+z x) \mathrm{d} x \mathrm{~d} y \mathrm{~d} z
\]

其中 \(V\) 是由 \(x \geq 0, y \geq 0,0 \leq z \leq 1\) 和 \(x^{2}+y^{2} \leq 1\) 所确定的空间区域．

\subsection{题目8}
九．（15 分）函数 \(f(x)\) 在 \((-\infty,+\infty)\) 二阶可导，\(f(x)\) 是 \((-\infty,+\infty)\) 上的有界函数，证明：存在\(\xi \in(-\infty,+\infty)\) ，满足 \(f^{\prime \prime}(\xi)=0\) ．

\subsection{题目9}
十．（5 分）函数 \(f(x)\) 在 \((-\infty,+\infty)\) 上各阶可导，且存在 \(M>0\) ，使得

\[
\left|f^{(k)}(x)\right| \leq M, \forall x \in(-\infty,+\infty), k=1,2, \cdots
\]

如果在一无限有界集 \(E\) 上，\(f(x) \equiv 0\) ，证明：在 \((-\infty,+\infty)\) 上，\(f(x) \equiv 0\) ．
\end{document}
